\documentclass[12pt,letterpaper]{article}

\usepackage[margin=1.25in]{geometry}
\usepackage{booktabs}
\usepackage{array}
\usepackage{setspace}

\usepackage[hidelinks]{hyperref}
\usepackage{natbib}
\usepackage{amsmath}

\onehalfspacing

\title{French and American Bimetallism in the Nineteenth Century:\\
Mint Ratios, Coinage Data, and the Road to 1873}

\author{}
\date{March 2026}

\begin{document}

\maketitle

\begin{abstract}
This note summarizes the operation of France's bimetallic currency system from 1803
to 1873, the legal and institutional framework that sustained it, and what is known
about the time-series of gold and silver coinage at the French mints.  It then
describes the parallel evolution of mint ratios at the United States Mint from 1792
to 1873, including the pivotal change enacted in the Coinage Act of 1834.  A final
section asks whether earlier and persistent coordination of the American and French
mint ratios at 15.5:1 would have produced a more robust bimetallic system.  The
answer the literature provides is nuanced: a US ratio permanently aligned with
France's 15.5:1 would have enlarged the bimetallic buffer in the critical decades of
the 1850s--1870s, but the secular rise in silver production after 1870 would
eventually have tested even an enlarged Franco-American bimetallic bloc.
\end{abstract}

\tableofcontents

\bigskip

%---------------------------------------------------------------------
\section{The Legal Foundation of French Bimetallism, 1803--1873}
%---------------------------------------------------------------------

\subsection{The Law of 7 Germinal An~XI (28 March 1803)}

France's bimetallic system rested on a single statutory instrument.  The Law of
7~Germinal An~XI, passed on 28~March 1803 under Napoleon as First Consul, defined
the franc Germinal in terms of both metals simultaneously:
\begin{align}
  1\text{ franc} &= 4.500 \text{ g fine silver}, \label{eq:silver}\\
  1\text{ franc} &= 0.2903 \text{ g fine gold}. \label{eq:gold}
\end{align}
Dividing \eqref{eq:silver} by \eqref{eq:gold} gives the statutory \emph{mint
ratio}---the number of grams of silver legally equivalent to one gram of gold:
\begin{equation}
  \rho_{\mathrm{FR}} = \frac{4.500}{0.2903} = 15.5. \label{eq:ratio}
\end{equation}
Equivalently, the law set the mint price of silver at 200 francs per kilogram
(at 90\% fineness) and of gold at 3{,}100 francs per kilogram (at 90\%
fineness)~\citep{wiegand2023}.  The ratio 15.5:1 had already been used for French
coinage since at least 1785~\citep{wikipedia_bimetallism} and was close to the
world market ratio in 1803.

\textbf{The ratio was never changed.}  From its enactment in 1803 until France
suspended free silver coinage in September 1873---a span of seventy years---the
statutory mint ratio remained fixed at 15.5:1~\citep{oppers2000,flandreau2004}.
No law adjusted it.  The entire burden of adjustment to changing world bullion
markets was borne by changes in the \emph{composition} of France's coin
circulation, not by changes in the price.

\subsection{Free Coinage and the Role of the Mints}

The law granted any member of the public the right to bring gold or silver bullion
to a government mint, have it assayed, and receive coins at the statutory price,
subject only to a small seigniorage charge to cover minting costs.  This
\emph{frappe libre} (free coinage) was the operational core of the system.  It
applied symmetrically to both metals.

The primary coinage facility was the Monnaie de Paris (Paris Mint), with provincial
mints in Bordeaux, Lyon, Strasbourg, and several other cities also active at
various times during the period.  The institutional division of responsibility
should be noted: the Bank of France, founded in 1800, held a monopoly on banknote
issue after 1848 and maintained large metallic reserves that backed its notes, but
the \emph{coinage terms}---who could bring metal, at what price, in what
denominations---were fixed by law and administered through the Ministry of Finance
and the Monnaie de Paris administration.  The Bank of France was not itself the
authority that set or enforced the mint ratio~\citep{ans2014}.

\citet{oppers2000} emphasizes the key mechanical point: the mints did not trade gold
for silver directly and did not attempt to fix the relative price of the two metals
through market operations.  They simply coined whatever bullion was presented to
them.  The relative price between gold and silver was sustained indirectly, through
the public's arbitrage between the mint and the bullion market.  If the market ratio
rose above 15.5, silver became cheap relative to the mint price of silver coins;
individuals would bring silver to the mint, expand the silver coin supply, and
export gold---exactly Gresham's law in action.  The converse applied when the market
ratio fell below 15.5.

\subsection{Denominations and the 1864 Partial Debasement}

The law of 1803 authorized gold coins in denominations of 20 and 40 francs (with 5,
10, 50, and 100 franc gold pieces added over subsequent decades), and silver coins of
5, 2, 1, and smaller denominations.  The workhorse gold coin was the 20-franc
Napoleon, first struck on 28~March 1803: weight 6.45161~g, fineness
0.900, hence containing 5.80645~g of fine gold.  Approximately 500 million 20-franc
gold coins were minted over the period 1803--1914~\citep{thomasnumismatics}.

A partial qualification of full bimetallism occurred in 1864, when all silver coins
\emph{except} the 5-franc piece were debased from 900\textperthousand{} to
835\textperthousand{} silver without changing their weights.  This converted
subsidiary silver (1F, 2F, 50~centimes, 20~centimes) into token coins whose face
value exceeded their bullion value; only the 5-franc coin remained a full-weight
commodity coin~\citep{wikipedia_frenchfranc}.  The formal 15.5 ratio continued to
apply to the 5-franc coin and to gold coinage.

\subsection{The Latin Monetary Union, 1865}

France extended its mint regime internationally by spearheading the Latin Monetary
Union (LMU), founded by treaty on 23~December 1865 among France, Belgium, Italy,
and Switzerland (Greece joined in 1867).  The LMU adopted the French specification
as the common standard: the same 15.5:1 bimetallic ratio, with the five-franc silver
coin at 900\textperthousand{} fineness and fractional silver at
835\textperthousand{}~\citep{wikipedia_lmu}.  The LMU took effect on
1~August~1866.

%---------------------------------------------------------------------
\section{Gresham's Law and the Three Phases of French Bimetallism}
%---------------------------------------------------------------------

The fixed 15.5 ratio interacted with a market ratio that varied over time.  The
theory is straightforward: whichever metal is undervalued at the mint (i.e., worth
more in bullion than the implied mint price) will be exported; whichever is
overvalued will be imported and coined in large quantities.  France in practice
cycled through three broad phases.

\paragraph{Phase 1 (1803--c.1848): Silver dominates; gold exported.}
From 1803 until around 1848, the world market ratio lay mostly above 15.5,
meaning gold was undervalued at the French mint relative to its bullion value.
French citizens exported gold and brought silver to the mint.  France was for most
of this period effectively on a \emph{de facto} silver standard, even though the
legal system permitted gold coinage.

\paragraph{Phase 2 (1848--c.1866): Gold dominates; silver exported.}
The California Gold Rush of 1848--1855 and the subsequent Australian discoveries
flooded world markets with gold and depressed its relative price.  The market ratio
fell below 15.5, making silver the dearer metal and triggering its export---
principally to India, which absorbed large quantities of silver throughout the 1850s
and 1860s.  France imported gold on an enormous scale, absorbing an estimated two-fifths of world
gold production between 1848 and 1870~\citep{wikipedia_bimetallism}.  By the 1850s
France was effectively on a gold standard in practice.

\paragraph{Phase 3 (1866--1873): Silver returns; gold begins to leave.}
As the initial wave of Californian and Australian production was absorbed, and as
global demand for silver (particularly in Asia) remained high, the market ratio
crept back above 15.5 after 1866.  Silver again began flowing into France and gold
flowed out---a reversal that intensified just as Germany was preparing its currency
reform.

Despite these large compositional swings, the market ratio remained remarkably
stable.  \citet{oppers2000} documents that over the fifty years from 1823 to 1873,
the world market ratio never departed more than about 6 percent from the French mint
ratio of 15.5, despite major fluctuations in the relative supplies of gold and
silver.  This remarkable stabilization is the central empirical puzzle addressed by
the target-zone models of \citet{oppers1995}, \citet{oppers2000}, and
\citet{velde_weber2000}.  The bimetallic mint ratio acted as a pair of barriers (an
upper and lower band) that the market ratio could not breach for long, because
profitable arbitrage would immediately be activated.

%---------------------------------------------------------------------
\section{Annual Coinage at the French Mints: Data Sources}
%---------------------------------------------------------------------

The mint ratio itself---15.5:1---was constant throughout 1803--1873, so there is no
time series on the \emph{ratio}.  What varied dramatically from year to year was how
much gold versus silver was coined.  Several quantitative snapshots are available
from secondary sources, but the full annual series require consulting the
specialized monographic literature.

\subsection{Known Quantitative Evidence}

The clearest quantitative illustration of the regime's dynamics comes from the
silver-inflow data around the break-point of 1873.  \citet{wikipedia_lmu} reports
that in all of 1871 and 1872 combined the French mint received only 5~million francs
worth of silver for coinage; in 1873 alone it received 154~million francs---a
roughly thirty-fold increase reflecting the collapse of the market ratio and the
looming threat that Germany would unload its silver onto the French bimetallic
buffer.

For the gold-inflow phase, the scale was also very large.  During the decade
1848--1858, French gold coinage ran at historically unprecedented rates as Californian
and Australian gold flooded in; France accumulated close to 40 percent of world gold
production over the 1848--1870 period as a whole~\citep{wikipedia_bimetallism}.

\subsection{Primary Data Sources in the Literature}

The authoritative sources for annual mint-output and specie-stock data are as
follows.

\begin{enumerate}
  \item \textbf{\citet{flandreau2004}}, Chapter~4, ``Coin Memories: New Estimates
  of France's Specie Stock (1840 to 1878),'' and the associated appendix tables.
  This is the most comprehensive reconstruction of annual French mint output and
  the decomposition of the specie stock into gold and silver, drawing on French
  archival material.  An earlier version of the same estimates appeared in
  \citet{flandreau1995}.

  \item \textbf{\citet{sicsic1989}}, ``Estimation du stock de monnaie m\'etallique
  en France \`a la fin du XIX\textsuperscript{e} si\`ecle,'' \textit{Revue
  \'economique}.  This is the standard reference for metallic money stocks in
  late-nineteenth-century France and is cited by virtually every quantitative
  paper in the literature.

  \item \textbf{\citet{velde_weber2000}}, whose data appendix reconstructs annual
  gold and silver coin stocks for France and other countries, drawing on
  \citet{sicsic1989}, \citet{flandreau1995}, and the annual reports of the
  US~Mint Director (which tracked foreign specie stocks in detail).

  \item \textbf{\citet{oppers1996}}, who constructed annual data on the composition
  of the French money supply (gold vs.\ silver shares) for the 1820s--1870s.

  \item \textbf{\citet{wiegand2022}}, Figure~7, which plots the share of silver
  coins in France's specie circulation from 1849 to 1873, offering an accessible
  graphical time series based on the Flandreau--Sicsic estimates.
\end{enumerate}

%---------------------------------------------------------------------
\section{The Collapse of Bimetallism, 1871--1878}
%---------------------------------------------------------------------

The end of the French bimetallic system came in a cascade of events triggered by the
Franco-Prussian War of 1870--71.

\begin{itemize}
  \item \textbf{July 1871.}  The Berlin Mint suspended silver coinage as Germany
  began accumulating gold using France's war indemnity payments.
  \item \textbf{July 1873.}  Germany formally adopted the gold standard
  (Reichstag legislation).
  \item \textbf{6~September 1873.}  France settled the final installment of its war
  indemnity, paying approximately 5 billion francs (exceeding 20 percent of French
  GDP)~\citep{wiegand2023}.  The following day the Paris Mint limited silver
  coinage, breaking the bimetallic bond.
  \item \textbf{30~January 1874.}  The LMU members agreed in Paris to limit free
  silver coinage temporarily.
  \item \textbf{1876.}  France suspended silver coinage entirely.
  \item \textbf{1878.}  The LMU formally suspended free silver coinage altogether,
  placing the union on a de facto gold standard~\citep{wikipedia_lmu}.
\end{itemize}

\citet{flandreau1996} argued in a widely cited paper that France's September 1873
decision was the critical event, not Germany's demonetization per se.  The
suspension ended France's function as the world's bimetallic stabilizer and
immediately caused silver to depreciate sharply against gold.  \citet{wiegand2022}
confirms that bimetallism was technically viable in September 1873---France's
absorption capacity was sufficient to handle the German silver---and that its
end reflected a strategic game between France and Germany rather than economic
inevitability.  \citet{meissner2015} reaches a less optimistic verdict, arguing
that by 1875 the growing number of countries shifting to gold would have made
bimetallism unsustainable regardless.

\bigskip

%---------------------------------------------------------------------
\section{The US Mint and Bimetallic Ratios, 1792--1873}
%---------------------------------------------------------------------

\subsection{The Coinage Act of 1792: Hamilton's 15:1 Ratio}

The United States adopted bimetallism from the outset under Alexander Hamilton's
monetary architecture.  Hamilton's \textit{Report on the Establishment of a Mint}
(January 1791) proposed free coinage of both gold and silver, with the dollar
defined in both metals simultaneously.  The Coinage Act of 2~April 1792 enacted
his plan, with Section~11 fixing the proportional value of gold to silver at
\emph{fifteen to one} by weight of pure metal~\citep{statutesandstories}.

Under this law the silver dollar contained 371.25 grains of pure silver (416 grains
standard weight), and the gold Eagle (\$10) contained 247.5 grains of pure gold
(270 grains standard weight).  The implied ratio was exactly 15:1.  Hamilton chose
this figure because it was close to the then-prevailing market ratio in the United
States and in continental Europe~\citep{econlib_silverperiod}.

The system immediately ran into Gresham's law.  The world market ratio in the
early nineteenth century was closer to 15.5:1, meaning gold was slightly
\emph{undervalued} at the American mint (15:1) relative to its European market
value.  Gold tended to be exported; the US operated on a de facto silver standard
for most of the period 1792--1834.  The 15:1 ratio was, from the outset, 0.5
percentage points below the French ratio of 15.5:1, and this small but persistent
divergence had large monetary consequences over the long run.

\subsection{The Coinage Act of 1834: Shift to 16:1}

Political pressure to bring gold back into domestic circulation led Congress to
overhaul the mint ratio in 1834.  The Coinage Act of 28~June 1834 reduced the gold
content of the gold Eagle by approximately 6.5 percent, raising the effective
mint ratio from 15:1 to approximately \emph{sixteen to one}~\citep{apmex_1834}.

The new ratio was not derived from the market ratio.  The market ratio at the time
was roughly 15.7:1, and proposals were made in Congress for a ratio of 15.6:1
(close to France's 15.5) or 15.625:1.  These were rejected; the ratio of 16:1 was
chosen to favor gold decisively and ensure that gold coins would circulate
domestically rather than be exported~\citep{econlib_1834}.  It was recognized by
contemporaries that 16:1 would drive silver out.  Mr.~Gorham of Massachusetts warned
on the floor of the House that making gold too cheap at the mint would cause silver
to disappear from circulation---precisely what happened~\citep{econlib_1834}.

The Coinage Act of 1837 made a minor technical adjustment, standardizing the
gold--silver alloy proportion at one-tenth for both metals, which altered the exact
ratio to approximately 15.988:1 (often rounded to 16:1 in the literature).

After 1834, with the market ratio typically in the range of 15.5--15.8, the
American mint ratio of 16:1 overvalued gold and undervalued silver.  Silver was
now the metal whose coin value fell below its bullion value; it tended to be melted
or exported.  The Coinage Act of 1853 addressed the resulting shortage of small
change by formally debasing fractional silver coins (dimes, quarters, half dollars)
to 835\textperthousand{} fineness---the same solution France had adopted in 1864.
Full-weight silver dollars were still legally authorised but were seldom coined.
By the 1850s the United States was on a de facto gold standard~\citep{legalclarity}.

\subsection{The Coinage Act of 1873: The ``Crime of '73''}

The Civil War suspended specie convertibility.  When Congress turned to restoring
it, the Coinage Act of 12~February 1873 omitted the standard silver dollar from
the list of authorized coins, effectively demonetizing silver and placing the
United States permanently on a gold standard upon resumption of specie payments
(which occurred in 1879).  The act was labeled the ``Crime of~'73'' by populists
and silver interests who argued it had been passed without adequate public debate
and that it inexcusably committed the country to the deflationary gold standard.

\citet{friedman1990jep} argued that had the United States retained bimetallism---
even at the 16:1 ratio rather than France's 15.5:1---resumption of specie payments
would likely have occurred on a silver rather than gold basis, and the severe
deflation of the 1870s--1890s would have been substantially moderated.  In the
Friedman counterfactual, the market gold-silver ratio would have remained close to
16:1 through roughly 1914, rather than rising to nearly 32:1 by 1896.

\subsection{Summary: US Mint Ratios}

Table~\ref{tab:us_ratios} summarizes the statutory mint ratios at the United States
Mint from 1792 through the abandonment of bimetallism.

\begin{table}[h!]
\centering
\caption{US Mint Ratios (Silver per Unit of Gold), 1792--1873}
\label{tab:us_ratios}
\begin{tabular}{lll}
\toprule
\textbf{Period} & \textbf{Statutory Ratio} & \textbf{Governing Legislation} \\
\midrule
1792--1834 & 15.00:1  & Coinage Act of 1792 \\
1834--1837 & $\approx$16.002:1 & Coinage Act of 1834 \\
1837--1873 & $\approx$15.988:1 & Coinage Act of 1837 (minor adjustment) \\
1873--     & Gold standard; silver dollar omitted & Coinage Act of 1873 \\
\bottomrule
\end{tabular}
\smallskip\\
\small\textit{Note}: The French mint ratio was 15.5:1 throughout 1803--1873.
\end{table}

\bigskip

%---------------------------------------------------------------------
\section{Would a Shared US--France Ratio of 15.5:1 Have Stabilized\\
International Bimetallism?}
%---------------------------------------------------------------------

The two legs of France's bimetallic system---the fixed mint ratio and free coinage
open to the public---gave France its role as the ``buyer and seller of last resort''
for both metals.  As \citet{friedman1990jep} put it, France was ``a major
participant in the market for silver and gold, important enough to be able to peg
the price ratio despite major changes in the relative production of silver and
gold.''  The United States, as the other large economy with a bimetallic legal
framework, could have played a complementary role---but its ratio diverged from
France's for virtually the entire period.

\subsection{The Three Gaps and Their Consequences}

Three distinct gaps between US and French ratios can be identified.

\textbf{1792--1803 (before French reform):} the US was at 15:1 while France was
also near 15.5:1 (the 1785 standard) before the 1803 codification.  The gap was
about 0.5 in ratio units.  At a market ratio near 15.5, gold was undervalued in
the US and tended to be exported; the US circulated silver, France circulated
silver; both behaved similarly.  The gap was small enough not to create dramatic
triangular arbitrage.

\textbf{1803--1834:} France at 15.5, US at 15.0.  With the market ratio near 15.5,
the US at 15:1 consistently undervalued gold, which flowed out.  Had the US been at
15.5, it would have been in the same regime as France: both would have faced
effective silver monometallism until the California gold rush.  The direction of
effect is not obviously destabilizing, but the US was not contributing to the
bimetallic buffer.

\textbf{1834--1873:} France at 15.5, US at $\approx$16.  This is the consequential
gap.  At a market ratio typically between 15.4 and 15.8 (before the post-1873
collapse), the US ratio of 16 consistently overvalued gold at the US mint and
undervalued silver.  As a result:
\begin{itemize}
  \item After 1848, when the market ratio fell below 15.5, gold flowed into
  France under the French rule but not into the US under the 16:1 rule (gold was
  overvalued at 16:1 even more, so gold was not exported from the US);
  \item The US effectively withdrew from the bimetallic buffer function.  Silver
  was scarce in US circulation even as it was abundant in France.
  \item The global bimetallic stabilization mechanism ran through France alone,
  making it more fragile.
\end{itemize}

\subsection{The Flandreau--Friedman--Meissner Counterfactual}

The literature provides a fairly direct answer to the counterfactual of a permanent
US ratio of 15.5:1.

\citet{friedman1990jep} showed, using a partial-equilibrium accounting framework,
that the US at 16:1 after 1834 was effectively a gold-standard country, since
silver was always departing from US circulation.  Retention of bimetallism at
15.5:1 throughout 1803--1873 would have meant the US shared France's buffer-stock
role.  The combined specie stock of France and the US was very large---together they
constituted the dominant bimetallic bloc.  A Franco-American bimetallic bloc at the
same ratio of 15.5:1 would have been approximately twice as large as the French
bloc alone in terms of monetary metal holdings.

\citet{meissner2015} provides a formal assessment using \citeauthor{flandreau1996}'s
(\citeyear{flandreau1996}) general equilibrium model.  The viability of bimetallism
depends on the size of the bimetallic bloc relative to the world supply of
precious metals.  A country defecting to gold raises the demand for gold and
lowers the demand for silver, which stresses the bimetallic anchor at the margin.
Meissner shows that a country only needs to have a specie demand equal to about
31 percent of France's to tip bimetallism from viable to unviable, and Germany's
specie demand in the early 1870s was approximately 27 percent of France's---very
close to that threshold.  Had the United States---with specie demand that by the
1880s exceeded 50 percent of France's---been part of a bimetallic bloc at 15.5:1,
the threshold for instability would have been much harder to reach.

\citet{friedman1990jep} concluded explicitly: ``The U.S.\ could have played the same
role after 1873 in stabilizing the gold-silver price ratio that France did before
1873.''  However, at the US ratio of 16:1 rather than 15.5:1, a coordinated bimetallic
system would have required a different ratio for both countries.  \citet{meissner2015}
notes that a joint Franco-American attempt to maintain bimetallism at 15.5:1 in 1878
would have been ``tenuous'' by that late date, given the number of countries already
committed to gold.

\citet{flandreau1996} and \citet{wiegand2022} identify two structural prerequisites
for stable international bimetallism: (i)~a large-enough bimetallic bloc, and
(ii)~a single shared mint ratio (since two different bimetallic ratios create
arbitrage possibilities that can rapidly drain one country of both metals).
A world in which the US had adopted France's 15.5:1 ratio in 1834---or even in
1792---would have satisfied both prerequisites far better than the actual world of
1834--1873 in which the US ratio lay a full half-percentage-point above France's.

\subsection{Caveats and the Long Run}

Several caveats temper an overly optimistic assessment.

First, even a large Franco-American bloc at 15.5:1 would eventually have confronted
the secular rise in silver production from the American West beginning in the late
1860s and accelerating through the 1880s.  \citet{velde_weber2000} show that the
long-run trajectory of the gold-silver ratio was determined by relative metal
production, and that by the 1880s the secular fall in silver's value would have
pushed a bimetallic equilibrium toward effective silver monometallism, eroding the
gold component of the money supply.  \citet{meissner2015} estimates this would
have become unviable by around 1875--1880 given realistic silver production
scenarios.

Second, the choice of a mint ratio is itself a political act, and the US political
economy of 1834 was hostile to any ratio that did not decisively favor gold
circulation.  The Jacksonian Democrats who passed the Coinage Act of 1834 were
motivated primarily by a desire to return gold to everyday use (replacing bank
notes) rather than by any theory of international monetary stabilization.
A counterfactual that asks Congress to have chosen 15.5:1 instead of 16:1 in 1834
must contend with this political economy.

Third, the Velde-Weber target-zone model shows that stability of the market ratio
requires not only a large bimetallic bloc but also that both metals actually
circulate simultaneously in that bloc---not just that a single metal dominates
most of the time with the other appearing only at the margin.  At 16:1, the US was
for decades in ``effective gold monometallism,'' meaning the US specie stock
contributed nothing to the bimetallic buffer even though the US was technically a
bimetallic country by statute.  At 15.5:1 with the US ratio matching France's, the
US would have been in the same effective regime as France---sometimes effective
silver monometallism, sometimes effective gold monometallism---but in either case
contributing additively to the buffer capacity.

\subsection{Assessment}

The scholarly consensus, synthesized from \citet{friedman1990jep},
\citet{flandreau1996,flandreau2004}, \citet{oppers1996,oppers2000},
\citet{velde_weber2000}, \citet{meissner2015}, and \citet{wiegand2022}, supports
the following conclusion.  Had the United States permanently adopted France's
ratio of 15.5:1 from 1792 (or at least from 1803), and maintained free coinage of
both metals throughout the pre-Civil War period:
\begin{enumerate}
  \item The bimetallic buffer bloc would have been substantially larger in the
  critical decade of the 1860s--1870s, making Germany's silver demonetization far
  less disruptive to the ratio;
  \item The US would likely have resumed specie convertibility on a silver rather
  than gold basis in 1876--1879, sustaining bimetallism for at least an additional
  decade;
  \item The global deflation of the 1880s--1890s---attributable in large part to
  the undersupply of gold relative to real economic growth under the gold
  standard---would have been significantly moderated;
  \item But bimetallism would still have faced severe stress from the silver glut
  of the 1880s--1890s, and its long-run viability beyond 1890 was uncertain even
  in the most optimistic counterfactual.
\end{enumerate}

In short, a coordinated US--France bimetallic standard at 15.5:1, maintained from
1803, would have been both more robust and more durable than the historical
outcome---but it would not have been indefinitely sustainable in the face of the
secular changes in precious metal supplies that unfolded after 1870.

\bigskip

%---------------------------------------------------------------------
\section{Chronological Summary}
%---------------------------------------------------------------------

Table~\ref{tab:timeline} summarizes the main events discussed in this note.

\begin{table}[h!]
\centering
\caption{Key Dates in French and American Bimetallism}
\label{tab:timeline}
\begin{tabular}{>{\raggedright}p{2.0cm}p{11.5cm}}
\toprule
\textbf{Date} & \textbf{Event} \\
\midrule
1785          & France adopts 15.5:1 mint ratio for its coinage. \\
2 Apr 1792    & Coinage Act of 1792: US sets ratio at 15:1 (Hamilton). \\
28 Mar 1803   & Law of 7 Germinal An~XI: France codifies 15.5:1; franc Germinal established. \\
1803--1848    & Market ratio $>$ 15.5: silver coined in France, gold exported; US effective silver standard. \\
1848--1866    & California/Australian gold rush: market ratio $<$ 15.5; gold coined in France, silver exported to Asia. \\
28 Jun 1834   & Coinage Act of 1834: US raises ratio to $\approx$16:1; US shifts to effective gold standard. \\
1837          & Coinage Act of 1837: minor alloy adjustment; ratio $\approx$15.988. \\
1853          & Coinage Act of 1853: US debases fractional silver coins. \\
1864          & France debases subsidiary silver (except 5-franc piece) from 900 to 835\textperthousand. \\
23 Dec 1865   & Treaty of the Latin Monetary Union: France, Belgium, Italy, Switzerland adopt 15.5:1. \\
1 Aug 1866    & Latin Monetary Union takes effect. \\
1866--1873    & Market ratio rises above 15.5 again; silver flows back to France. \\
12 Feb 1873   & US Coinage Act of 1873 (``Crime of '73''): silver dollar omitted; US on gold. \\
Jul 1873      & Germany formally adopts gold standard. \\
6 Sep 1873    & Paris Mint limits silver coinage; bimetallic bond broken. \\
30 Jan 1874   & LMU limits free silver coinage. \\
1876          & France suspends silver coinage entirely. \\
1878          & LMU formally adopts gold standard. \\
\bottomrule
\end{tabular}
\end{table}

\bigskip

\bibliographystyle{plainnat}
\bibliography{bimetallism}

\end{document}
